% SGA 3, Expose V, Loop 1 English reconstruction.
% Scope: proof of Theorem 4.1(iii), then the opening of part f) through
% Lemma 4.1.3 and its proof; stops before the proof of Theorem 4.1(iv).
% Authority: Polo--Gille Exp5-13oct24.pdf, component-local pp. 15--16
% (combined-reader pp. 277--278). Re-edited page headers 265--266;
% internal original-page marker 267 occurs on local p. 15.
% Jacob Reinhold's CC BY 4.0 Markdown is comparison/drafting material only.

现在证明 \hyperref[thm:sga3-V-4.1-iii]{定理~4.1(iii)}.
再次回顾一下, \(\underline P\) 表示概形 \(P\) 的底集,
\(\underline d:\underline P\to\underline Q\) 为态射 \(d:P\to Q\)
诱导的映射. \hyperref[proof:sga3-V-4.1-b]{b)} 可重述为:
分量为 \(\underline{d_0}\) 和 \(\underline{d_1}\) 的映射
\[
  \underline{d_0}\boxtimes\underline{d_1}:
  \underline{X_1}\longrightarrow
  \underline{X_0}\times_{\underline Y}\underline{X_0}
\]
是满射. 此映射分解为
\[
  \underline{X_1}
  \xrightarrow{\ \underline{d_0}\boxtimes\underline{d_1}\ }
  \underline{X_0}\times_{\underline Y}\underline{X_0}
  \xrightarrow{\ q\ }
  \underline{X_0\times_YX_0},
\]
其中 \(q\) 是典范映射. 因而 \(\underline{d_0}\boxtimes\underline{d_1}\)
的像包含每个满足
\[
  \{v\}=q^{-1}(q(v)).
\]
的点 \(v\in\underline{X_0\times_YX_0}\).
特别地, 当 \(v\) \emph{在 \(Y\) 上有理},
即剩余域 \(\kappa(v)\) 可等同于 \(v\) 在 \(Y\) 中的像 \(w\)
的剩余域 \(\kappa(w)\) 时, 此条件成立.%
\footnote{\label{note:sga3-V-23}编者注: Lecture Notes in Mathematics
151 的页面在这里排错了顺序; 实际顺序为
265--266--268--269--267--270--271.}

若 \(v\in\underline{X_0\times_YX_0}\) 在 \(Y\) 上不有理,
仍用 \(w\) 表示它在 \(Y\) 中的像.
由引理~\hyperref[lem:sga3-V-4.1.2]{4.1.2},
存在根基为 \(\mathfrak m\) 的局部环 \(C\) 及局部平坦同态
\[
  f:\mathcal O_{Y,w}\longrightarrow C
\]
使 \(C/\mathfrak m\) 作为 \(\kappa(w)\)-代数同构于 \(\kappa(v)\).
令 \(Y'=\operatorname{Spec}C\), 并令 \(\pi:Y'\to Y\) 为 \(f\) 诱导的态射.
典范投影
\[
  (X_0\times_YX_0)\times_YY'\longrightarrow X_0\times_YX_0
\]
将某个在 \(Y'\) 上有理的点 \(v'\) 映到 \(v\).
由于
\[
  (X_0\times_YX_0)\times_YY'
  \simeq
  (X_0\times_YY')\times_{Y'}(X_0\times_YY'),
\]
又由于 \hyperref[thm:sga3-V-4.1]{定理~4.1} 的假设及前面的结果,
特别是 \hyperref[proof:sga3-V-4.1-b]{b)},
在基变换 \(\pi:Y'\to Y\) 后仍成立,
点 \(v'\) 是某个元素
\[
  u'\in\underline{X_1\times_YY'}
\]
在 \(d_0\boxtimes d_1\) 的基变换下的像.
若 \(u\) 是 \(u'\) 在 \(X_1\) 中的像, 则
\[
  v=(d_0\boxtimes d_1)(u),
\]
这就是所需结论.

\medskip
\phantomsection\label{proof:sga3-V-4.1-f}%
\noindent\textbf{f) \hyperref[thm:sga3-V-4.1-iv]{(iv)} 的证明.}%
\footnote{\label{note:sga3-V-24}编者注:
\hyperref[lem:sga3-V-4.1.3]{下面的引理} 是从
\hyperref[bib:sga3-V-DG70]{[DG70]}, I, \S5, 命题~1.5 补入的;
另见 EGA \(\mathrm{IV}_3\), 8.11.5 的证明.
原文隐含地使用了它, 而 \hyperref[bib:sga3-V-DG70]{[DG70]}, III,
\S2, 4.6 中明确使用了它.
它在 \hyperref[thm:sga3-V-7.1]{下面的定理~7.1} 中也有用.}

\medskip
\phantomsection\label{lem:sga3-V-4.1.3}%
\noindent\textbf{引理 4.1.3.}---
\textit{若概形的单态射 \(f:T\to Z\) 有限, 则它是闭浸入.}

用仿射开子集 \(Z_i\) 覆盖 \(Z\), 并将 \(f\) 替换为诱导态射
\(f^{-1}(Z_i)\to Z_i\), 由于有限态射仿射,
即可化归到以下情形:
\[
  Z=\operatorname{Spec}B,\qquad T=\operatorname{Spec}A.
\]
由于 \(f\) 是单态射, 对角态射
\[
  T\longrightarrow T\times_ZT
\]
为同构 (EGA~I, 5.3.8); 等价地, 乘法给出一个同构
\[
  A\otimes_BA\xrightarrow{\ \sim\ }A.
\]
因此, 对 \(B\) 的每个极大理想 \(\mathfrak m\), 有同构
\[
  (A/\mathfrak mA)\otimes_k(A/\mathfrak mA)
  \xrightarrow{\ \sim\ }A/\mathfrak mA,
  \qquad k=B/\mathfrak m.
\]
由于 \(A\) 在 \(B\) 上有限, \(k\)-向量空间 \(A/\mathfrak mA\)
具有有限维数 \(d\), 而此同构蕴含 \(d=0\) 或 \(d=1\).
因而映射
\[
  k=B/\mathfrak m\longrightarrow A/\mathfrak mA
\]
满射. 在 \(\mathfrak m\) 处局部化后应用 Nakayama 引理,
可知 \(B_{\mathfrak m}\to A_{\mathfrak m}\) 满射.
由此得到 \(B\to A\) 满射: 其余核 \(C\) 对每个极大理想
\(\mathfrak m\) 都满足 \(C_{\mathfrak m}=0\), 因而为零.
这就证明了 \hyperref[lem:sga3-V-4.1.3]{本引理}.
% Translation statistics: source_lines=123; output_lines=122; inline_math=59; display_math=13; equation=0; equation_star=0; align=0; footnotes=2.
% Review: complete assigned component; pending_TODO=0; all math symbols, labels, references, diagram paths and environments checked; formula prose translated.
